\section{常曲率流形}\label{sec:manifolds-of-constant-curvature}

回顾, 若一个流形在所有点的所有截面曲率都取同一个常数值 $K$, 则称它为常曲率流形. 利用曲率形式可以容易地判别这些流形.

\begin{RiemannLemma}\label{lem:curvature-forms-characterize-constant-curvature}
若 $M$ 是曲率为 $K$ 的常曲率流形, 则在每点 $p\in M$ 附近, 所有曲率形式 $\Omega_i^j$ 都满足
\begin{equation}\label{eq:curvature-forms-in-constant-curvature}
  \Omega_i^j=-K\,\omega^i\wedge\omega^j,
\end{equation}
其中 $\{\omega^1,\ldots,\omega^n\}$ 是定义在 $p$ 的一个邻域上的任意标准正交余标架场. 反之, 若 $M$ 的每点附近都有一个标准正交标架场 $\{E_1,\ldots,E_n\}$, 相应的余标架场 $\{\omega^1,\ldots,\omega^n\}$ 对某个常数 $K$ 满足 \eqref{eq:curvature-forms-in-constant-curvature}, 则 $M$ 的曲率恒为 $K$.
\end{RiemannLemma}

\begin{RiemannProof}
若 $M$ 的曲率恒为 $K$, 则
\begin{align*}
  \Omega_i^j&=\sum_{k<l}\Omega_i^j(E_k,E_l)\,\omega^k\wedge\omega^l
  =\sum_{k<l}\omega^j(R(E_k,E_l)E_i)\,\omega^k\wedge\omega^l\\
  &=\sum_{k<l}\langle R(E_k,E_l)E_i,E_j\rangle\,\omega^k\wedge\omega^l
  =\sum_{k<l}R_{klij}\,\omega^k\wedge\omega^l\\
  &=-\sum_{k<l}K(\delta_{ki}\delta_{lj}-\delta_{kj}\delta_{li})\,\omega^k\wedge\omega^l
  =-K\,\omega^i\wedge\omega^j.
\end{align*}
反之, 设存在常数 $K$, 使在每点 $p\in M$ 的某个邻域上有 $\Omega_i^j=-K\,\omega^i\wedge\omega^j$. 则对切空间 $T_pM$ 的每个截面 $\Pi$, 相应的截面曲率由下式给出:
\[
  K(\Pi)=-R(X,Y,X,Y)
\]
其中 $X,Y$ 是张成 $\Pi$ 的两个线性无关向量 (假设它们张成单位面积的平行四边形). 用 $p$ 附近的标准正交标架场, 有 $X=\sum_{i=1}^nX^iE_i$, $Y=\sum_{i=1}^nY^iE_i$, 所以
\begin{align*}
  K(\Pi)&=-\sum_{i,j,k,l=1}^n X^iY^jX^kY^lR(E_i,E_j,E_k,E_l)\\
  &=-\sum_{i,j,k,l=1}^n X^iY^jX^kY^l\Omega_k^l(E_i,E_j)\\
  &=K\sum_{i,j,k,l=1}^n X^iY^jX^kY^l\omega^k\wedge\omega^l(E_i,E_j)\\
  &=K\sum_{i,j,k,l=1}^n X^iY^jX^kY^l
    \bigl(\omega^k(E_i)\omega^l(E_j)-\omega^k(E_j)\omega^l(E_i)\bigr)\\
  &=K\sum_{i,j,k,l=1}^n X^iY^jX^kY^l(\delta_{ik}\delta_{jl}-\delta_{jk}\delta_{il})\\
  &=K\bigl(\lVert X\rVert^2\lVert Y\rVert^2-\langle X,Y\rangle^2\bigr)=K.
\end{align*}
\end{RiemannProof}

现在来看一个应用此引理的例子.

\begin{RiemannExample}\label{ex:constant-curvature-of-hyperbolic-space}
半径为 $a>0$ 的 $n$ 维\textbf{双曲空间} $H^n(a)$ 是赋有 Riemann 度量
\[
  g_{ij}(x)=\frac{a^2}{(x^n)^2}\delta_{ij}
\]
的开半空间
\[
  \{(x^1,\ldots,x^n)\in\mathbb{R}^n\mid x^n>0\}.
\]
这个 Riemann 流形的截面曲率恒为 $K=-\frac1{a^2}$. 事实上, 利用上述引理, 将证明 $H^n(a)$ 上存在标准正交标架场 $\{E_1,\ldots,E_n\}$, 其对偶余标架场 $\{\omega^1,\ldots,\omega^n\}$ 对 $K=-\frac1{a^2}$ 满足
\begin{equation}\label{eq:target-curvature-forms-of-hyperbolic-space}
  \Omega_i^j=-K\,\omega^i\wedge\omega^j.
\end{equation}
为此, 考虑自然坐标系 $x:H^n(a)\to\mathbb{R}^n$, 及由 $X_i=\frac{\partial}{\partial x^i}$ 构成的相应坐标标架场 $\{X_1,\ldots,X_n\}$. 由于
\[
  \langle X_i,X_j\rangle=\frac{a^2}{(x^n)^2}\delta_{ij},
\]
得到一个标准正交标架场 $\{E_1,\ldots,E_n\}$, 其中 $E_i=\frac{x^n}{a}X_i$, 而相应的对偶余标架场 $\{\omega^1,\ldots,\omega^n\}$ 满足 $\omega^i=\frac{a}{x^n}\,\mathrm{d}x^i$. 则
\[
  \mathrm{d}\omega^i=\frac{a}{(x^n)^2}\,\mathrm{d}x^i\wedge\mathrm{d}x^n
  =\frac1a\,\omega^i\wedge\omega^n
  =\sum_{j=1}^n\omega^j\wedge\left(-\frac1a\delta_{jn}\omega^i\right),
\]
所以, 利用结构方程
\begin{align*}
  \mathrm{d}\omega^i&=\sum_{j=1}^n\omega^j\wedge\omega_j^i\\
  \omega_i^j&+\omega_j^i=0,
\end{align*}
可以猜测联络形式由 $\omega_j^i=\frac1a(\delta_{in}\omega^j-\delta_{jn}\omega^i)$ 给出. 事实上, 容易验证这些形式满足上述结构方程, 因而由方程解的唯一性, 它们必为联络形式. 有了这些形式, 现在容易用第三个结构方程计算曲率形式 $\Omega_i^j$:
\[
  \mathrm{d}\omega_i^j=\sum_{k=1}^n\omega_i^k\wedge\omega_k^j+\Omega_i^j.
\]
有
\[
  \mathrm{d}\omega_i^j=\mathrm{d}\left(\frac1a(\delta_{jn}\omega^i-\delta_{in}\omega^j)\right)
  =\frac1{a^2}(\delta_{jn}\,\omega^i\wedge\omega^n-\delta_{in}\,\omega^j\wedge\omega^n)
\]
和
\begin{align*}
  \sum_{k=1}^n\omega_i^k\wedge\omega_k^j
  &=\frac1{a^2}\sum_{k=1}^n(\delta_{kn}\omega^i-\delta_{in}\omega^k)
    \wedge(\delta_{jn}\omega^k-\delta_{kn}\omega^j)\\
  &=\frac1{a^2}\sum_{k=1}^n(\delta_{kn}\delta_{jn}\,\omega^i\wedge\omega^k
    -\delta_{kn}\,\omega^i\wedge\omega^j+\delta_{in}\delta_{kn}\,\omega^k\wedge\omega^j)\\
  &=\frac1{a^2}(\delta_{jn}\,\omega^i\wedge\omega^n-\omega^i\wedge\omega^j
    +\delta_{in}\,\omega^n\wedge\omega^j),
\end{align*}
因而
\[
  \Omega_i^j=\mathrm{d}\omega_i^j-\sum_{k=1}^n\omega_i^k\wedge\omega_k^j
  =\frac1{a^2}\,\omega^i\wedge\omega^j.
\]
由此得到 $K=-\frac1{a^2}$.
\end{RiemannExample}

Euclidean 空间 $\mathbb{R}^n$ 的曲率恒为零. 此外, 容易看出半径为 $r$ 的球面 $S^n(r)\subset\mathbb{R}^{n+1}$ 的曲率恒为 $\frac1{r^2}$ [参见 Exercise~\RGref{ex:isometric-immersions}(\RGref{ex:constant-curvature-of-round-sphere})]. 因此在任意维数下, 都有曲率为任意负常数 ($H^n(a)$), 零 ($\mathbb{R}^n$) 或正的常数 ($S^n(r)$) 的流形例子. 注意, 这些例子全都单连通且测地完备. 事实上, Euclidean 空间 $\mathbb{R}^n$ 的测地线的像是直线, $S^n(r)$ 是紧的, 而 $H^n(a)$ 的测地线的像或者是垂直于平面 $x^n=0$ 且圆心在该平面上的半圆, 或者是从平面 $x^n=0$ 出发的竖直半直线 [参见 Exercise~\ref{ex:manifolds-of-constant-curvature}(\ref{ex:geodesics-of-hyperbolic-space})].

每个单连通且测地完备的常曲率流形都与上述一个例子等距, 如下述定理所示. 一般而言, 若该流形不单连通 (但仍测地完备), 则它与上述某个例子在其等距映射群的一个离散子群的自由固有作用下的商空间等距 (可以证明 Riemann 流形的等距映射群总是 Lie 群).

\begin{RiemannTheorem}[Killing-Hopf]\label{thm:killing-hopf}
设 $M$ 是曲率恒为 $K$ 的连通且测地完备的 $n$ 维 Riemann 流形.
\begin{enumerate}[label=(\arabic*),leftmargin=*]
\item 若 $M$ 单连通, 则它与下列流形之一等距: 当 $K>0$ 时为 $S^n\bigl(\frac1{\sqrt K}\bigr)$, 当 $K=0$ 时为 $\mathbb{R}^n$, 当 $K<0$ 时为 $H^n\bigl(\frac1{\sqrt{-K}}\bigr)$.
\item 若 $M$ 不单连通, 则 $M$ 与商空间 $\widetilde M/\Gamma$ 等距, 其中 $\widetilde M$ 是上述单连通流形之一, $\Gamma$ 是 $\widetilde M$ 的等距映射群的非平凡离散子群, 在 $\widetilde M$ 上固有且自由地作用.
\end{enumerate}
\end{RiemannTheorem}

\begin{RiemannProof}
本定理的证明见 \RGcite{ch04-dC93}. 这里只给出 $M$ 单连通, $n=2$ 且 $K=0$ 情形下的证明. 在此情形下, Cartan–Hadamard 定理 [参见第~\ref{chap:riemannian-manifolds}~章的 Exercise~\ref{ex:hopf-rinow-theorem}(\ref{ex:ambrose-and-cartan-hadamard-theorems})] 蕴含, 给定 $p\in M$, 映射 $\exp_p:T_pM\to M$ 是微分同胚. 设 $\{E_1,E_2\}$ 是 $M$ 上的一个全局标准正交标架 (由全局 Cartesian 坐标的相应标架正交归一化得到). 由于 $K=0$, 相应的联络形式 $\omega_1^2$ 满足 $\mathrm{d}\omega_1^2=0$, 因而由 Poincaré 引理 [参见第~\ref{chap:differential-forms}~章的 Exercise~\ref{ex:differential-forms}(\ref{ex:de-rham-cohomology})], 对某个光滑函数 $f\in C^\infty(M)$, 有 $\omega_1^2=\mathrm{d}f$. 设 $\{F_1,F_2\}$ 是与 $\{E_1,E_2\}$ 同定向的标准正交标架, 使得 $E_1$ 与 $F_1$ 之间的夹角为 $\theta=-f$. 则其联络形式 $\overline\omega{}_1^2$ 满足 $\overline\omega{}_1^2=\omega_1^2+\mathrm{d}\theta=0$, 即 $\nabla_{F_1}F_1=\nabla_{F_1}F_2=\nabla_{F_2}F_1=\nabla_{F_2}F_2=0$. 从而 $[F_1,F_2]=0$, 因此由第~\ref{chap:differentiable-manifolds}~章的 Theorem~\ref{thm:commuting-flows-and-lie-brackets}, 它们的流可交换. 于是, 可以用参数化
\[
  \varphi(x,y)=\psi_{1,x}\circ\psi_{2,x}(p)
\]
在 $M$ 上引入局部坐标 $(x,y)$ (其中 $\psi_1,\psi_2$ 是 $F_1,F_2$ 的流). 利用 $\frac{\partial}{\partial x}=F_1$ 和 $\frac{\partial}{\partial y}=F_2$, 容易证明 $\varphi(x,y)=\exp_p\bigl(x(F_1)_p+y(F_2)_p\bigr)$, 所以 $(x,y)$ 实际上是全局坐标. 由于在这些坐标下度量写成
\[
  g=\mathrm{d}x\otimes\mathrm{d}x+\mathrm{d}y\otimes\mathrm{d}y,
\]
因此 $M$ 与 $\mathbb{R}^2$ 等距.
\end{RiemannProof}

\begin{RiemannExample}\label{ex:complete-flat-surfaces}
设 $\widetilde M=\mathbb{R}^2$. 则等距映射的子群 $\Gamma$ 不能含有带不动点的等距映射 (因为它自由作用). 因此 $\Gamma$ 只能包含平移和滑移反射 (即反射后再沿反射轴的方向平移). 此外, $\Gamma$ 至多由两个元素生成, 其中一个可以假定为平移 [参见 Exercise~\ref{ex:manifolds-of-constant-curvature}(\ref{ex:discrete-isometry-groups-of-plane})]. 因此:
\begin{enumerate}[label=(\arabic*),leftmargin=*]
\item 若 $\Gamma$ 由一次平移生成, 则所得曲面为圆柱面;
\item 若 $\Gamma$ 由两次平移生成, 则得到环面;
\item 若 $\Gamma$ 由一次滑移反射生成, 则得到 Möbius 带;
\item 若 $\Gamma$ 由一次平移和一次滑移反射生成, 则得到 Klein 瓶.
\end{enumerate}
这些就是测地完备 Euclidean 曲面 (曲率恒为零的 2 维流形) 的所有可能例子.
\end{RiemannExample}

\begin{RiemannExample}\label{ex:orientable-hyperbolic-surfaces}
双曲平面 $H^2$ 的保持定向的等距映射群是 $\mathit{PSL}(2,\mathbb{R})=\mathit{SL}(2,\mathbb{R})/\{\pm\operatorname{id}\}$, 它在 $H^2$ 上的作用为
\[
  \begin{pmatrix}a&b\\c&d\end{pmatrix}\cdot z:=\frac{az+b}{cz+d},
\]
其中将 $\mathbb{R}^2\cong\mathbb{C}$ 等同 [参见 Exercise~\ref{ex:manifolds-of-constant-curvature}(\ref{ex:orientation-preserving-isometries-of-hyperbolic-plane}) 及第~\RGref{subsec:notes-manifolds-of-constant-curvature}~节]. 要找出可定向双曲曲面, 即曲率恒为 $K=-1$ 的曲面, 必须找到 $\mathit{PSL}(2,\mathbb{R})$ 在 $H^2$ 上固有且自由地作用的离散子群 $\Gamma$. 这里的可能性要多得多. 例如, 可以考虑由平移 $f(z)=z+2\pi$ 生成的群 $\Gamma=\langle f\rangle$. 所得曲面称为\textbf{伪球面}, 与圆柱面同胚 (参见 Fig.~\ref{fig:pseudosphere}). 然而, $y\to+\infty$ 的一端宽度趋于零, 而 $y\to0$ 的一端宽度趋于 $+\infty$. 它朝向两端的高度都是无限的. 注意, 此曲面上存在横截自交有限次的测地线 (参见 Fig.~\ref{fig:pseudosphere-geodesics}).
\end{RiemannExample}
\RGfigure[.73\textwidth]{pseudosphere}{伪球面}
\RGfigure[.65\textwidth]{pseudosphere-geodesics}{伪球面上的测地线轨迹}

取双曲多边形 (以测地线为边界), 并用等距映射把它们的边等同, 可得到其他例子. 例如, Fig.~\ref{fig:hyperbolic-polygon-punctured-sphere}(b) 的曲面是由等距映射 $g(z)=z+2$ 和 $h(z)=\frac{z}{2z+1}$ 等同 Fig.~\ref{fig:hyperbolic-polygon-punctured-sphere}(a) 中多边形的边得到的. 选取其他多边形可以得到紧双曲曲面. 事实上, 对任意 Euler 特征数为负的拓扑 2 维流形, 都存在与它同胚的紧双曲曲面 (此时 Gauss–Bonnet 定理不允许非负的 Euler 特征数).
\RGfigure[.73\textwidth]{hyperbolic-polygon-punctured-sphere}{(a) 双曲多边形; (b) 三穿孔球面}

\begin{RiemannExample}\label{ex:even-dimensional-positive-curvature-quotients}
要找出常正曲率的 Riemann 流形, 必须找到球面的等距映射群中固有且自由地作用的离散子群. 考虑 $K=1$ 的情形. 此时 $\Gamma\subset O(n+1)$ [参见 Exercise~\ref{ex:manifolds-of-constant-curvature}(\ref{ex:isometry-group-of-sphere})]. 因为它必须在 $S^n$ 上自由作用, 所以 $\Gamma\setminus\{\operatorname{id}\}$ 中的元素都不能有特征值 $1$. 将看到, 若 $n$ 为偶数, $S^n$ 和 $\mathbb{R}P^n$ 是仅有的曲率恒为 $1$ 的测地完备流形. 事实上, 若 $A\in\Gamma$, 则 $A$ 是一个 $(n+1)\times(n+1)$ 正交矩阵, 故其所有特征值的绝对值均为 $1$. 此外, 它的特征多项式的次数 $n+1$ 为奇数, 因而有一个等于 $\pm1$ 的实根. 因此 $A^2$ 有特征值 $1$, 所以 $A^2$ 必为恒等映射. 于是 $A=A^{-1}=A^t$, 因而 $A$ 对称, 从而所有特征值均为实数. 那么 $A$ 的特征值要么全等于 $1$ (若 $A=\operatorname{id}$), 要么全等于 $-1$, 此时 $A=-\operatorname{id}$. 因而 $\Gamma=\{\pm\operatorname{id}\}$, 这蕴含所得流形是 $S^n$ 或 $\mathbb{R}P^n$. 若 $n$ 为奇数, 还有其他可能, 其分类见 \RGcite{Wol78}.
\end{RiemannExample}

\begin{RiemannExercise}\label{ex:manifolds-of-constant-curvature}
\begin{enumerate}[label=(\arabic*),ref=\arabic*,leftmargin=*]
\item\label{ex:hyperbolic-metric-as-left-invariant-metric} 将 $(x^1,\ldots,x^n)\in H^n(a)$ 与由
\[
  g(t^1,\ldots,t^{n-1})=x^n(t^1,\ldots,t^{n-1})+(x^1,\ldots,x^{n-1})
\]
给出的仿射映射 $g:\mathbb{R}^{n-1}\to\mathbb{R}^{n-1}$ 等同, 可诱导出 Lie 群结构. 证明 $H^n(a)$ 的度量对该 Lie 群结构是左不变度量.
\item\label{ex:connection-forms-for-common-one-form} 证明若标准正交余标架场中的形式 $\omega^i$ 满足 $\mathrm{d}\omega^i=\alpha\wedge\omega^i$ (其中 $\alpha$ 是 1-形式), 则联络形式 $\omega_i^j$ 由 $\omega_i^j=\alpha(E_i)\omega^j-\alpha(E_j)\omega^i=-\omega_j^i$ 给出. 用此结论验证 Example~\ref{ex:constant-curvature-of-hyperbolic-space} 的结果.
\item\label{ex:constant-curvature-of-conformal-euclidean-metric} 设 $K$ 是实数, 令 $\rho=1+(\frac K4)\sum_{i=1}^n(x^i)^2$. 证明 $\mathbb{R}^n$ 上由
\[
  g_{ij}(p)=\frac1{\rho^2}\delta_{ij}
\]
定义的 Riemann 度量的截面曲率恒为 $K$.
\item\label{ex:geodesics-of-hyperbolic-space} 证明保持坐标函数 $x^n$ 不变的 Euclidean 空间 $\mathbb{R}^n$ 的任意等距映射都是 $H^n(a)$ 的等距映射. 用此结论确定 $H^n(a)$ 的所有测地线.
\item\label{ex:schur-theorem} (\emph{Schur 定理}) 设 $M$ 是维数 $n\geq3$ 的连通各向同性 Riemann 流形. 证明 $M$ 为常曲率流形.
{\small (\textbf{提示:} 利用结构方程证明 $\mathrm{d}K=0$).}
\item\label{ex:discrete-isometry-groups-of-plane} 为补全 Example~\ref{ex:complete-flat-surfaces} 中的细节, 证明:
\begin{enumerate}[label=(\alph*),leftmargin=*]
\item $\mathbb{R}^2$ 中没有不动点的等距映射或者是平移, 或者是滑移反射;
\item 在 $\mathbb{R}^2$ 上固有且自由地作用的任何离散等距映射群至多由两个元素生成, 其中一个可以假定为平移.
\end{enumerate}
\item\label{ex:mobius-band-and-klein-bottle-quotients} 设 $f,g:\mathbb{R}^2\to\mathbb{R}^2$ 是等距映射
\[
  f(x,y)=(-x,y+1)\quad\text{和}\quad g(x,y)=(x+1,y)
\]
(故 $f$ 是滑移反射, $g$ 是平移). 核实 $\mathbb{R}^2/\langle f\rangle$ 与 Möbius 带 (无边界) 同胚, $\mathbb{R}^2/\langle f,g\rangle$ 与 Klein 瓶同胚.
\item\label{ex:orientation-preserving-isometries-of-hyperbolic-plane} 设 $H^2$ 是双曲平面. 证明:
\begin{enumerate}[label=(\alph*),leftmargin=*]
\item 公式
\[
  \begin{pmatrix}a&b\\c&d\end{pmatrix}\cdot z:=\frac{az+b}{cz+d}\qquad(ad-bc=1)
\]
定义了 $\mathit{PSL}(2,\mathbb{R}):=\mathit{SL}(2,\mathbb{R})/\{\pm\operatorname{id}\}$ 在 $H^2$ 上通过保持定向的等距映射给出的作用;
\item 对任意两条以弧长为参数的测地线 $c_1,c_2:\mathbb{R}\to H^2$, 存在 $g\in\mathit{PSL}(2,\mathbb{R})$, 使得对所有 $s\in\mathbb{R}$ 有 $c_1(s)=g\cdot c_2(s)$;
\item 给定满足 $d(z_1,z_2)=d(z_3,z_4)$ 的 $z_1,z_2,z_3,z_4\in H^2$, 存在 $g\in\mathit{PSL}(2,\mathbb{R})$, 使得 $g\cdot z_1=z_3$ 且 $g\cdot z_2=z_4$;
\item 若 $H^2$ 的一个保持定向的等距映射有两个不动点, 则它必为恒等映射. 由此推出, 所有保持定向的等距映射均为 $f(z)=g\cdot z$ 的形式, 其中 $g\in\mathit{PSL}(2,\mathbb{R})$.
\end{enumerate}
\item\label{ex:identifications-on-hyperbolic-polygon} 核实 Example~\ref{ex:orientable-hyperbolic-surfaces} 中双曲平面的等距映射 $g(z)=z+2$ 和 $h(z)=\frac{z}{2z+1}$ 将 Fig.~\ref{fig:hyperbolic-polygon-punctured-sphere} 中的双曲多边形的边等同.
\item\label{ex:tractrix-and-tractroid} \textbf{曳物线}是由以下参数方程描述的曲线:
\[
  \begin{cases}
    x=u-\tanh u\\
    y=\operatorname{sech}u
  \end{cases}
  \qquad(u>0)
\]
(其名称源于如下性质: 曲线上任意一点沿切线到 $x$ 轴的距离恒为 $1$). 证明将曳物线绕 $x$ 轴旋转生成的旋转曲面 (\textbf{曳物旋转面}) 的 Gauss 曲率恒为 $K=-1$. 确定伪球面中与曳物旋转面等距的一个开子集.
{\small (\textbf{Remark:} 曳物旋转面并不测地完备. 事实上, Hilbert 于 1901 年证明嵌入 Euclidean 3 维空间的任何常负曲率曲面必不完备).}
\item\label{ex:isometry-group-of-sphere} 证明 $S^n$ 的等距映射群是 $O(n+1)$.
\item\label{ex:compact-two-dimensional-lie-groups} 设 $G$ 是 2 维紧 Lie 群. 证明:
\begin{enumerate}[label=(\alph*),leftmargin=*]
\item $G$ 可定向;
\item $\chi(G)=0$;
\item $G$ 上的任意左不变度量都是常曲率的;
\item $G$ 是 2 维环面 $T^2$.
\end{enumerate}
\end{enumerate}
\end{RiemannExercise}
